\documentclass[11pt]{article}
\usepackage[T1]{fontenc}
\usepackage{lmodern}
\usepackage{amsmath,amssymb,amsthm,mathtools}
\usepackage[margin=1in]{geometry}
\usepackage{hyperref}
\newtheorem{theorem}{Theorem}[section]
\newtheorem{proposition}[theorem]{Proposition}
\newtheorem{lemma}[theorem]{Lemma}
\newtheorem{corollary}[theorem]{Corollary}
\theoremstyle{definition}
\newtheorem{definition}[theorem]{Definition}
\newtheorem{remark}[theorem]{Remark}
\newcommand{\HH}{\mathsf H}
\newcommand{\HHZ}{\mathsf H_0}
\newcommand{\SSS}{\mathsf{CSRig}}
\newcommand{\CR}{\mathsf{CRing}}
\newcommand{\Om}{\Omega}
\newcommand{\one}{\mathbf 1}
\newcommand{\Con}{\operatorname{Con}}
\newcommand{\Fix}{\operatorname{Fix}}
\newcommand{\Ann}{\operatorname{Ann}}
\title{The literal third absorber: exact adjunctions, receivers,\newline congruences, and a distinct zero-preserving variant}
\author{}
\date{19 September 2026}
\begin{document}
\maketitle

\noindent\textbf{Scope and source.}
This is a standalone independent derivation of the third-element construction
in Section 9 of \emph{An Examination of Semiring Structures Derived from the
Integers by Sequential Adjunction of Identity and Absorbing Elements},
Zenodo version record 17547186, recovered source file
\texttt{17547186\_5.tex}. That source explicitly calls its literal extension a
commutative unital hemiring and explicitly states that it is not a standard
semiring. Its defining equation is \(\tau\Om=\Om\), not \(\tau\Om=\tau\).
The proof below retains that equation. A construction with the latter equation
appears only in the final section and is a different object.
All rings in the main construction are commutative and unital. We allow the
zero ring when describing quotients; statements needing \(0_R\ne1_R\) say so.

\section{Precisely specified categories and the original two-level object}

\begin{definition}\label{def:categories}
An object of \(\HH\) is a set \(D\) with two binary operations, written
\(+\) and multiplication, and constants \(0_D,1_D\), such that
\((D,+,0_D)\) and \((D,\cdot,1_D)\) are commutative monoids and
\(a(b+c)=ab+ac\) for all \(a,b,c\in D\).
The equation \(0_Da=0_D\) is \emph{not} an axiom of \(\HH\).
A morphism in \(\HH\) preserves both binary operations and both constants.
The category \(\HHZ\) has the same objects, with maps preserving both binary
operations and \(0_D\), but not necessarily \(1_D\).
The category \(\SSS\) is the full subcategory of \(\HH\) satisfying
\(0_Da=0_D\) for every \(a\in D\); these are the commutative unital
standard semirings in this note. The singleton algebra is allowed.
The category \(\CR\) has commutative unital rings and unital ring maps.
\end{definition}

These explicit axioms fix the meaning of the weakened category regardless of
terminological differences concerning the word ``hemiring.''
Only binary-operation-preserving maps will be called such when either constant
is not required to be preserved.

\begin{definition}\label{def:G}
Let \(R\) be a commutative unital ring, put \(e=0_R\), and choose
\(\tau\notin R\). The disjoint union \(G(R)=R\sqcup\{\tau\}\) has the
original ring operations on every pair of elements of \(R\), and has
\[
 \tau+x=x+\tau=x,\qquad \tau x=x\tau=\tau\quad(x\in G(R)).
\]
Its proposed additive identity is \(\tau\), and its proposed multiplicative
identity is \(1_R\). In particular \(e\ne\tau\), including when \(R\) is
the zero ring.
\end{definition}

\begin{proposition}\label{prop:G}
The displayed operations make \(G(R)\) an object of \(\SSS\). Furthermore,
for \(a,b\in G(R)\),
\[
 a+b=\tau\ \Longleftrightarrow\ a=b=\tau,
 \qquad
 ab=\tau\ \Longleftrightarrow\ (a=\tau\text{ or }b=\tau).
\tag{1}\label{eq:Gsupport}
\]
These equivalences hold even if \(R\) has zero divisors or nonzero elements
annihilated by positive integers.
\end{proposition}
\begin{proof}
Commutativity follows from commutativity in \(R\) and the symmetric rules
involving \(\tau\). In an additive triple, deleting occurrences of
\(\tau\) leaves the same ordered sum of ring elements under either
parenthesization; a triple consisting entirely of \(\tau\) has value
\(\tau\). Thus addition is associative and has identity \(\tau\).
In a multiplicative triple with an occurrence of \(\tau\), both
parenthesizations have value \(\tau\). A triple in \(R\) is associative
by the ring axiom. The element \(1_R\) acts as identity on \(R\) and on
\(\tau\).

For distributivity, if \(a=\tau\), both \(a(b+c)\) and \(ab+ac\) are
\(\tau\). Suppose \(a\in R\). If \(b,c\in R\), distributivity is the
original ring equation. If \(b=\tau,c\in R\), its two sides are
\(ac\) and \(\tau+ac=ac\). If \(b\in R,c=\tau\), its two sides are
\(ab\) and \(ab+\tau=ab\). If \(b=c=\tau\), its two sides are
\(\tau\) and \(\tau+\tau=\tau\). These cases exhaust all triples.
The required absorbing-zero axiom is one of the defining rules.

A sum with at least one entry in \(R\) lies in \(R\), whereas
\(\tau+\tau=\tau\). A product of two entries in \(R\) lies in \(R\),
whereas every product having an entry \(\tau\) is \(\tau\).
This proves both equivalences. A ring product equal to \(e\) is still in
\(R\), and does not equal \(\tau\).
\end{proof}

\begin{proposition}\label{prop:ringprojection}
The map
\[
 p_R:G(R)\longrightarrow R,\qquad p_R(\tau)=e,\quad p_R(r)=r\ (r\in R)
\tag{2}\label{eq:ringprojection}
\]
is a surjective morphism in \(\SSS\). Its kernel congruence has the one
nonsingleton class \(\{\tau,e\}\). The set inclusion \(R\hookrightarrow
G(R)\) preserves both binary operations and the multiplicative identity,
but does not preserve the additive identity.
\end{proposition}
\begin{proof}
On a pair from \(R\), the map is the identity, hence preserves both
operations. On \((\tau,r)\), the additive equation is \(r=e+r\),
and the multiplicative equation is \(e=er\). On \((\tau,\tau)\),
they are \(e=e+e\) and \(e=e^2\). It sends the source zero \(\tau\)
to the target zero \(e\), and sends \(1_R\) to \(1_R\).
It is surjective because it fixes every ring element. Its fibers are
\(\{\tau,e\}\) and the other singleton ring elements. The claims about
the inclusion follow directly from the restricted operations and
\(e\ne\tau\).
\end{proof}

\section{The literal adjunction of an absorber for both operations}

\begin{definition}\label{def:A}
For an arbitrary \(D\in\HH\), let \(A(D)=D\sqcup\{\Om\}\), where
\(\Om\) is a new element. Retain both operations on \(D\), and set
\[
 \Om+x=x+\Om=\Om,\qquad \Om x=x\Om=\Om\quad(x\in A(D)).
\tag{3}\label{eq:literal}
\]
Retain the constants \(0_D,1_D\). For a ring \(R\), write
\[
 F(R)=A(G(R))=R\sqcup\{\tau,\Om\}.
\]
Thus in \(F(R)\), \(\tau\) remains the additive identity,
\(1_R\) remains the multiplicative identity, and
\(\tau\Om=\Om\tau=\Om\).
\end{definition}

\begin{theorem}\label{thm:Aaxioms}
For every \(D\in\HH\), the algebra \(A(D)\) is an object of \(\HH\),
and the inclusion \(\eta_D:D\hookrightarrow A(D)\) is an injective
\(\HH\)-morphism. The new element \(\Om\) is the unique global absorber
for either operation. The object \(A(D)\) is never an object of \(\SSS\).
In particular the literal source construction \(F(R)\) is consistent
for every commutative unital ring \(R\).
\end{theorem}
\begin{proof}
Both extended operations are commutative. In a triple containing \(\Om\),
both parenthesizations of either operation are \(\Om\); a triple wholly
in \(D\) satisfies associativity there. The constant \(0_D\) remains
an additive identity because \(0_D+\Om=\Om\), and \(1_D\) remains a
multiplicative identity because \(1_D\Om=\Om\).

For distributivity, take \(a,b,c\in A(D)\). If all three are in \(D\),
use distributivity in \(D\). If \(a=\Om\), the left side is \(\Om\)
and the right side is \(\Om+\Om=\Om\). If \(a\ne\Om\) and
\(b=\Om\), the left side is \(a\Om=\Om\), and the right side is
\(\Om+ac=\Om\). If \(a,b\ne\Om,c=\Om\), the left side is
\(a\Om=\Om\), and the right side is \(ab+\Om=\Om\). This covers
all cases; right distributivity follows from commutative multiplication.
The inclusion preserves both operations and constants by construction.

If \(z\) were any additive absorber, then \(z=z+\Om=\Om\).
If \(z\) were any multiplicative absorber, then \(z=z\Om=\Om\).
Finally \(0_D\Om=\Om\ne0_D\), so global zero absorption fails.
No equation already holding between elements of \(D\) was changed.
\end{proof}

\begin{definition}
Let \(\HH_{\mathrm{abs}}\) have objects \((B,b)\), where \(B\in\HH\)
and \(b\in B\) satisfies \(b+y=b\) and \(by=b\) for every \(y\in B\).
Morphisms are \(\HH\)-morphisms preserving the distinguished absorber.
Let \(U:\HH_{\mathrm{abs}}\to\HH\) forget that designation.
\end{definition}

\begin{theorem}[Exact free-adjunction property]\label{thm:Auniversal}
For every \(D\in\HH\) and \((B,b)\in\HH_{\mathrm{abs}}\), restriction
along \(\eta_D\) gives a bijection
\[
 \operatorname{Hom}_{\HH_{\mathrm{abs}}}((A(D),\Om),(B,b))
 \ \longrightarrow\ \operatorname{Hom}_{\HH}(D,B).
\tag{4}\label{eq:adjunction}
\]
Its inverse sends \(f:D\to B\) to the unique function
\(\widehat f:A(D)\to B\) with \(\widehat f|_D=f\) and
\(\widehat f(\Om)=b\). These bijections commute with precomposition
and postcomposition; consequently \(A\) is left adjoint to \(U\).
\end{theorem}
\begin{proof}
The function \(\widehat f\) is uniquely determined because its domain is
the disjoint union \(D\sqcup\{\Om\}\). For a pair in \(D\), preservation
of either operation is exactly the corresponding equation for \(f\).
For a pair with an occurrence of \(\Om\), the image of the source result
is \(b\); the operation on the images is also \(b\), by its absorbing
property. This includes the pair \((\Om,\Om)\). The constants lie in
\(D\) and are preserved by \(f\). Hence \(\widehat f\) is a morphism.
Restriction and extension are inverse as functions on the indicated hom-sets.
For \(g:D'\to D\), both \(\widehat f\circ A(g)\) and
\(\widehat{f\circ g}\) equal \(f\circ g\) on \(D'\) and send the new
absorber to \(b\); they coincide. For an absorber-preserving
\(h:(B,b)\to(C,c)\), both \(h\circ\widehat f\) and
\(\widehat{h\circ f}\) restrict to \(h\circ f\) and send the new
absorber to \(c\); they coincide. These are the required naturality
equations.
\end{proof}

\begin{proposition}[Receivers without a globally designated absorber]
Let \(f:D\to B\) be an \(\HH\)-morphism and let \(b\in B\). There
exists an \(\HH\)-morphism \(h:A(D)\to B\) restricting to \(f\) and
sending \(\Om\) to \(b\) if and only if
\[
 b+b=b,\quad b^2=b,\quad b+f(x)=b,\quad bf(x)=b\quad(x\in D).
\tag{5}\label{eq:relative-receiver}
\]
When it exists it is unique.
\end{proposition}
\begin{proof}
Necessity follows by applying \(h\) to each defining equation involving
\(\Om\), including its pair with itself. For sufficiency, define \(h\)
on the disjoint union as specified. Pairs from \(D\) are handled by
\(f\), mixed pairs by the last two displayed equations, and the pair
\((\Om,\Om)\) by the first two. Constants are in \(D\), so they
are preserved. There is no other possible value at any domain element.
\end{proof}

\section{Exact receivers in standard semirings and rings}

\begin{theorem}[Standard-semiring reflection]\label{thm:reflection}
Let \(D\in\HH\) and \(B\in\SSS\). Every \(\HHZ\)-morphism
\(h:A(D)\to B\) is the constant map with value \(0_B\).
That constant map is indeed an \(\HHZ\)-morphism. An \(\HH\)-morphism
\(A(D)\to B\) exists exactly when \(B\) is the singleton algebra;
then it is unique. Thus the universal standard-semiring receiver of
\(A(D)\), for unital zero-preserving morphisms, is the singleton semiring.
The universal ring receiver is also the zero ring.
\end{theorem}
\begin{proof}
Since \(\Om0_D=\Om\), preservation of multiplication and the additive
identity gives
\[
 h(\Om)=h(\Om)h(0_D)=h(\Om)0_B=0_B.
\]
For every \(x\in A(D)\), the equation \(\Om+x=\Om\) then gives
\(0_B+h(x)=0_B\), whence \(h(x)=0_B\). Conversely the constant-zero
map preserves addition, multiplication, and zero in a standard semiring.
If the map also preserves the multiplicative identity, it forces
\(1_B=0_B\). For each \(y\in B\), the standard-semiring laws yield
\(y=1_By=0_By=0_B\), so \(B\) is the singleton algebra. The unique
map to a singleton preserves all required operations and constants.

For the asserted universal property, denote the singleton semiring by
\(\one\). The unique map \(A(D)\to\one\) factors every unital
zero-preserving map to a standard semiring: its target has just been proved
to be a singleton, and the factor \(\one\to B\) is unique. This is
the universal receiver, not merely an obstruction to an injective map.
Rings are standard semirings, so the same argument gives the ring assertion.
\end{proof}

\begin{proposition}[Even maps to rings forgetting the constants collapse]
If \(K\) is a ring, every map \(h:A(D)\to K\) preserving the two
binary operations is the constant map \(0_K\), whether or not either
constant was required to be preserved.
\end{proposition}
\begin{proof}
The equation \(\Om+x=\Om\) implies
\(h(\Om)+h(x)=h(\Om)\). Addition in \(K\) is a group; adding
\(-h(\Om)\) to both sides gives \(h(x)=0_K\) for every \(x\).
The constant-zero map preserves both binary operations in a ring.
\end{proof}

\begin{corollary}[What fails if both absorption axioms are demanded]
Suppose one asks for a standard semiring containing the original \(G(R)\)
with its additive zero \(\tau\), and asks simultaneously for a new element
\(\Om\) satisfying \(\Om\tau=\Om\) and \(\Om+x=\Om\) for every
element. The universal algebra satisfying those requirements is the
singleton semiring, and there is no such injective extension of \(G(R)\).
This conclusion does not invalidate the original \(\HH\)-construction.
\end{corollary}
\begin{proof}
In a standard semiring, \(\Om\tau=\tau\); the other demanded equation
gives \(\Om=\tau\). For every \(x\), addition then gives
\(x=\tau+x=\Om+x=\Om=\tau\). The singleton algebra satisfies all
the equations. It is their universal receiver by
Theorem~\ref{thm:reflection}. The original set \(G(R)\) contains the
distinct elements \(e,\tau\), so cannot inject into that receiver.
The original \(\HH\)-construction omits the premise of global zero
absorption and has already been verified directly.
\end{proof}

\section{Ring information is retained exactly in the weakened category}

\begin{theorem}[Exact quotient identifying the two original zeros]
\label{thm:quotientA}
There is a surjective \(\HH\)-morphism
\[
 A(p_R):F(R)\longrightarrow A(R),\qquad
 \tau\longmapsto e,\quad r\longmapsto r\ (r\in R),\quad
 \Om\longmapsto\Om.
\tag{6}\label{eq:Aprojection}
\]
Its kernel is precisely the congruence generated by \(\tau=e\), and
\[
 F(R)/(\tau=e)\ \cong\ A(R)
\tag{7}\label{eq:Aquotient}
\]
in \(\HH\). In particular this quotient has not collapsed the new
absorber. For any \(B\in\HH\), an \(\HH\)-morphism \(h:F(R)\to B\)
factors through \(A(p_R)\) exactly when \(h(\tau)=h(e)\); its factor
is unique.
\end{theorem}
\begin{proof}
The map \(p_R\) is a morphism by Proposition~\ref{prop:ringprojection}.
Extending it by the new absorber, as in Theorem~\ref{thm:Auniversal},
gives the displayed morphism. Its fibers are \(\{\tau,e\}\), all
other singleton ring elements, and the singleton \(\{\Om\}\).
Thus its kernel contains \((\tau,e)\) and no other nonequal pair
except its symmetric pair. Every equivalence relation containing
\((\tau,e)\) contains this kernel, so this is exactly the least
congruence containing the specified equation.

If \(h(\tau)=h(e)\), define its factor on \(A(R)\) by taking the common
value on each fiber. To verify preservation of an operation, choose
preimages of its two arguments under \(A(p_R)\); the operation of the
preimages maps to the operation of the arguments because \(A(p_R)\)
is a homomorphism, and applying \(h\) gives the required equality.
The same reasoning for constants shows that the factor is in \(\HH\).
Surjectivity makes the factor unique. Conversely any factor is constant
on the fiber \(\{\tau,e\}\). The induced bijection from the quotient
to \(A(R)\) preserves operations and constants, as does its inverse by
the same fiber description.
\end{proof}

\begin{theorem}[Full faithfulness]\label{thm:fullyfaithful}
For every pair of commutative unital rings \(R,R'\), including zero rings,
the assignment
\[
 h\longmapsto F(h),\qquad
 F(h)(\tau)=\tau',\quad F(h)(r)=h(r),\quad F(h)(\Om)=\Om'
\tag{8}\label{eq:Ffunctor}
\]
is a bijection
\[
 \operatorname{Hom}_{\CR}(R,R')
 \ \cong\ \operatorname{Hom}_{\HH}(F(R),F(R')).
\tag{9}\label{eq:fullyfaithful}
\]
It preserves identities and compositions. Thus \(F\) is a fully faithful
functor, despite its singleton standard-semiring reflection.
\end{theorem}
\begin{proof}
For a ring homomorphism \(h\), the displayed extension preserves operations
on pairs of ring elements by the ring-homomorphism equations. A pair
having an occurrence of \(\Om\) maps on both sides to \(\Om'\).
A pair having \(\tau\) and no \(\Om\) is handled by the additive
identity and multiplicative-absorption rules in \(G(R')\).
Both constants are preserved. Hence \(F(h)\) is an \(\HH\)-morphism.

Conversely let \(f:F(R)\to F(R')\) be an \(\HH\)-morphism. Write
\(e'=0_{R'}\). It sends \(\tau\) to \(\tau'\) and \(1_R\) to
\(1_{R'}\). The equation \(\Om\tau=\Om\) implies
\(f(\Om)\tau'=f(\Om)\). In \(F(R')\), multiplication by \(\tau'\)
sends every element of \(G(R')\) to \(\tau'\) and sends \(\Om'\)
to \(\Om'\). Hence \(f(\Om)\) is either \(\tau'\) or \(\Om'\).
The first possibility, together with \(\Om+1_R=\Om\), would imply
\(\tau'+1_{R'}=\tau'\), contradicting \(1_{R'}\in R'\) and
\(\tau'\notin R'\). Therefore \(f(\Om)=\Om'\).

The additive idempotents of \(F(R')\) are exactly
\(\tau',e',\Om'\): for a ring element \(r'\), the equation
\(r'+r'=r'\) implies \(r'=e'\) by the ring additive inverse.
Since \(e+e=e\), the element \(f(e)\) is in this three-element set.
It cannot equal \(\Om'\), because \(e\tau=\tau\) would give
\(\Om'\tau'=\tau'\). If \(f(e)=\tau'\), then applying \(f\)
to \(1_R+(-1_R)=e\) gives a sum equal to \(\tau'\) with an entry
\(1_{R'}\in R'\). Such a sum cannot be \(\tau'\), by the defining
addition rules. Consequently \(f(e)=e'\).

For \(r\in R\), the equation \(er=e\) gives \(e'f(r)=e'\).
The value \(f(r)=\tau'\) would make its left side \(\tau'\), and
the value \(f(r)=\Om'\) would make it \(\Om'\). Both are impossible.
Thus \(f(r)\in R'\) for every \(r\in R\). The restriction
\(h=f|_R:R\to R'\) preserves addition, multiplication, zero, and unit;
it is a unital ring homomorphism. Its extension agrees with \(f\) on
each element of the disjoint union. This proves existence and uniqueness
of the inverse correspondence. The formulas show directly that
\(F(\mathrm{id}_R)=\mathrm{id}_{F(R)}\) and
\(F(kh)=F(k)F(h)\).
\end{proof}

\section{Complete congruence classification of the literal extension}

A congruence here is an equivalence relation compatible with both binary
operations. Quotient constants are the classes of the original constants.
An ideal of a ring contains its ring zero, is an additive subgroup, and
is closed under multiplication by all ring elements.

\begin{theorem}\label{thm:congruences}
For every ideal \(I\) of \(R\), including \(I=R\), define congruences
\(C_I,D_I\) of \(F(R)\) as follows:
\begin{itemize}
\item \(C_I\) has singleton classes \(\{\tau\}\) and \(\{\Om\}\),
and its classes in \(R\) are the additive cosets of \(I\).
\item \(D_I\) has singleton class \(\{\Om\}\), the class
\(\{\tau\}\cup I\), and the other additive cosets of \(I\) in \(R\).
\end{itemize}
Every congruence of \(F(R)\) is exactly one of these, or the universal
congruence. The quotient identifications are
\[
 F(R)/C_I\cong F(R/I),\qquad F(R)/D_I\cong A(R/I).
\tag{10}\label{eq:congquotients}
\]
The complete order among proper congruences is
\[
\begin{aligned}
 C_I\subseteq C_J&\Longleftrightarrow I\subseteq J,&
 D_I\subseteq D_J&\Longleftrightarrow I\subseteq J,\\
 C_I\subseteq D_J&\Longleftrightarrow I\subseteq J,&
 D_I\subseteq C_J&\text{ never holds.}
\end{aligned}
\tag{11}\label{eq:congorder}
\]
Above all of these is the universal congruence.
\end{theorem}
\begin{proof}
Let \(\rho\) be a congruence and suppose \(\Om\mathrel\rho x\) for
some \(x\in G(R)\). Multiplication by \(\tau\) gives
\(\Om\mathrel\rho\tau\), because \(x\tau=\tau\) in \(G(R)\).
For any \(y\in F(R)\), addition of \(y\) then gives
\(\Om=\Om+y\mathrel\rho\tau+y=y\). Thus \(\rho\) is universal.
Every proper congruence therefore has \(\Om\) as a singleton class.
Its restriction \(\sigma\) to \(G(R)\) is a congruence there. Conversely,
every congruence on \(G(R)\), even its universal one, extends to a
proper congruence on \(F(R)\) by taking \(\Om\) as a singleton: operations
with an occurrence of \(\Om\) give \(\Om\) on both sides of each
compatibility equation.

Restrict \(\sigma\) further to pairs in \(R\), and put
\(I=\{r\in R:r\mathrel\sigma e\}\). This set contains \(e\).
If \(a,b\in I\), then \(a+b\mathrel\sigma e+e=e\).
If \(a\in I\), multiplication by the fixed ring element \(-1_R\)
gives \(-a\mathrel\sigma e\). For any \(r\in R,a\in I\),
\(ra\mathrel\sigma re=e\). Hence \(I\) is a ring ideal. For
\(r,s\in R\), if \(r\mathrel\sigma s\), adding \(-s\) gives
\(r-s\mathrel\sigma e\), so \(r-s\in I\). Conversely, if
\(r-s\mathrel\sigma e\), adding \(s\) gives
\(r\mathrel\sigma s\). Thus the classes within \(R\) are precisely
the cosets of \(I\).

If \(\tau\) is a singleton under \(\sigma\), this is exactly the
restriction of \(C_I\). Otherwise \(\tau\mathrel\sigma r\) for
some \(r\in R\). Multiplying by \(e\) gives
\(\tau\mathrel\sigma e\). Therefore for every \(s\in R\),
\(s\mathrel\sigma\tau\) holds exactly when
\(s\mathrel\sigma e\), namely exactly when \(s\in I\).
This is exactly the restriction of \(D_I\), with no other possible
cross-stratum equivalences.

To verify existence and the quotient descriptions, the maps
\[
 \begin{array}{rclcrcl}
 q_I:F(R)&\longrightarrow&F(R/I),&&
 v_I:F(R)&\longrightarrow&A(R/I),\\
 r&\longmapsto&r+I,&&r&\longmapsto&r+I,\\
 \tau&\longmapsto&\tau_I,&&\tau&\longmapsto&I,\\
 \Om&\longmapsto&\Om_I,&&\Om&\longmapsto&\Om_I
 \end{array}
\]
are surjective \(\HH\)-morphisms. The first is \(F\) applied to the
ring quotient map, and its preservation equations also follow directly
from \eqref{eq:Ffunctor}. The second is the extension by the absorber
of the standard-semiring map \(G(R)\to R/I\) given by
\(r\mapsto r+I,\tau\mapsto I\); its verification is exactly the
ring-projection calculation with coefficients taken modulo \(I\).
Their fibers are the asserted classes, proving that \(C_I,D_I\) are
congruences and establishing the quotient isomorphisms. Their kernels
determine \(I\) by restriction to the ring-zero class; also \(C_I\)
and \(D_J\) are distinct because only the latter identifies \(\tau,e\).

The first three order equivalences follow by inspecting ring cosets:
refinement implies that every \(i\in I\), equivalent to \(e\),
lies in \(J\), and \(I\subseteq J\) gives the required refinement
of every displayed class. Finally \(D_I\) always identifies
\(\tau,e\), while no \(C_J\) does. This proves the last assertion
and completes the classification.
\end{proof}

\begin{corollary}
Forcing \(\tau=e\) in \(F(R)\) gives the proper quotient \(A(R)\),
whereas forcing \(\tau=\Om\), or forcing \(e=\Om\), gives the
singleton quotient. Requiring global absorption of the additive identity
therefore generates the universal congruence.
\end{corollary}
\begin{proof}
The first congruence is \(D_{\{e\}}\) by the theorem. Either of the
other equations identifies \(\Om\) with an element of \(G(R)\),
which the proof shows generates the universal congruence. Global
absorption of \(\tau\) includes \(\tau\Om=\tau\); the existing
equation \(\tau\Om=\Om\) then gives \(\tau=\Om\).
\end{proof}

\section{The triple of singlets, with its actual identity maps}

\begin{proposition}\label{prop:unitsfixed}
The group of multiplicative units of \(F(R)\) is exactly \(R^\times\).
Its multiplication action on \(F(R)\) is by automorphisms of the
additive monoid. The fixed set of the element \(-1_R\) is
\[
 \Fix(-1_R)=\{\tau,\Om\}\sqcup R[2],\qquad
 R[2]=\{r\in R:2r=e\}.
\tag{12}\label{eq:signfixed}
\]
The fixed set of the whole unit group is
\[
 \{\tau,\Om\}\sqcup
 \{r\in R:(u-1_R)r=e\text{ for every }u\in R^\times\}.
\tag{13}\label{eq:allfixed}
\]
In particular for \(R=\mathbb Z\) the unit group is \(\{1,-1\}\)
and the fixed set is exactly \(\{\tau,e,\Om\}\).
\end{proposition}
\begin{proof}
No product having a factor \(\Om\) is \(1_R\). A product having a
factor \(\tau\) and no factor \(\Om\) is \(\tau\ne1_R\).
Consequently both factors of a product equal to \(1_R\) must be in
\(R\); they are precisely inverse units there. For \(u\in R^\times\),
the map \(x\mapsto ux\) preserves addition by distributivity, fixes
\(\tau\), and has inverse \(x\mapsto u^{-1}x\) by associativity
and the identity axiom. The group-action composition equation is
\(u(vx)=(uv)x\).

Both \(\tau\) and \(\Om\) are fixed by every such \(u\), directly
from the multiplication rules. For \(r\in R\), multiplication by
\(-1_R\) gives \(-r\); equality \(-r=r\) is equivalent, by adding
\(r\), to \(2r=e\). Likewise \(ur=r\) is equivalent to
\((u-1_R)r=e\). This proves the two formulas.
For an integer unit \(n\), an integer \(m\) satisfies \(nm=1\);
taking ordinary absolute values gives \(|n||m|=1\), so \(n=1\) or
\(n=-1\), and both are units. The only integer satisfying \(2n=0\)
is \(n=0\), proving the final assertion.
\end{proof}

\begin{theorem}\label{thm:tripleretract}
For every \(R\), the subset \(E=\{\tau,e,\Om\}\) is closed under
both operations, with the following complete tables:
\[
\begin{array}{c|ccc}
 +&\tau&e&\Om\\\hline
 \tau&\tau&e&\Om\\
 e&e&e&\Om\\
 \Om&\Om&\Om&\Om
\end{array}
\qquad
\begin{array}{c|ccc}
 \cdot&\tau&e&\Om\\\hline
 \tau&\tau&\tau&\Om\\
 e&\tau&e&\Om\\
 \Om&\Om&\Om&\Om
\end{array}.
\tag{14}\label{eq:Etable}
\]
As an object of \(\HH\), it has \(0_E=\tau\) and \(1_E=e\).
Its additive order, defined by \(a\le b\) exactly when \(a+b=b\),
is the chain \(\tau<e<\Om\). There are typed maps
\[
 d_R:F(R)\longrightarrow E,\quad
 d_R(\tau)=\tau,\ d_R(r)=e,\ d_R(\Om)=\Om,
 \qquad j_R:E\hookrightarrow F(R).
\tag{15}\label{eq:dretract}
\]
Here \(d_R\) is an \(\HH\)-morphism; \(j_R\) is an
\(\HHZ\)-morphism; and \(d_Rj_R=\mathrm{id}_E\).
If \(R\) is nonzero, \(j_R\) is not unital, and no unital morphism
\(E\to F(R)\) exists at all. Thus this is a split retract in
\(\HHZ\), not in \(\HH\), when \(R\) is nonzero.
\end{theorem}
\begin{proof}
All entries involving \(\Om\) follow from its two absorbing laws.
The remaining entries follow from the definition of \(G(R)\) and
\(e+e=e^2=e\). These give exactly the displayed tables and closure.
The associative, commutative, and distributive equations restrict from
\(F(R)\). The tables exhibit \(\tau\) as the additive identity and
\(e\) as the multiplicative identity on this subset. They also exhibit
the stated three-element order, with no reverse inequalities between
distinct elements.

On the ambient set, multiplication by \(e\) sends \(\tau\) to
\(\tau\), every \(r\in R\) to \(e\), and \(\Om\) to \(\Om\).
It is therefore exactly \(j_Rd_R\). Distributivity gives
\(e(x+y)=ex+ey\). Since \(e^2=e\),
\((ex)(ey)=e^2xy=e(xy)\). It sends the ambient zero to \(\tau\),
and it sends \(1_R\) to \(e=1_E\). These equations prove that
\(d_R\) preserves operations and both constants with codomain \(E\).
The inclusion preserves the inherited binary operations and the common
additive identity. The composite \(d_Rj_R\) fixes each of the three
listed elements, so is the identity.

When \(R\) is nonzero, \(j_R(1_E)=e\ne1_R\). More strongly, a unital
map \(s:E\to F(R)\) would send \(e\) to \(1_R\); applying it to
\(e+e=e\) would give \(1_R+1_R=1_R\). This is an equation between
ring elements and implies \(1_R=e\) by adding \(-1_R\) in \(R\),
contrary to the hypothesis. Thus there cannot be a unital section or
even a unital map of the stated type.
\end{proof}

\begin{remark}
For a general ring, the triple \(E\) is always fixed by all units, but
it need not be the full fixed set. Formula~\eqref{eq:signfixed} retains
all possible two-torsion. For instance in characteristic two, \(-1_R=1_R\)
and every ring element is fixed by that particular element. Nothing in
the preceding proof discards such elements or changes the ring.
\end{remark}

\section{A different construction: preserve the global multiplicative zero}

This section does \emph{not} alter the literal object just proved.
It analyzes the alternative requirement that the original additive identity
remain a global multiplicative absorber after a third element is adjoined.

\begin{definition}\label{def:sharp}
Let \(K\in\SSS\), with \(0_K\ne1_K\). On
\(K^\sharp=K\sqcup\{w\}\), preserve the operations on \(K\), set
\(w+x=x+w=w\) for every \(x\in K^\sharp\), and define
\[
 0_Kw=w0_K=0_K,\qquad
 wa=aw=w\ (a\in K\setminus\{0_K\}),\qquad w^2=w.
\tag{16}\label{eq:sharp}
\]
The proposed constants are \(0_K,1_K\). A standard semiring \(K\)
is called zerosumfree here when \(a+b=0_K\) implies
\(a=b=0_K\); it has no nonzero zero divisors here when \(ab=0_K\)
implies \(a=0_K\) or \(b=0_K\).
\end{definition}

\begin{theorem}[Necessary and sufficient conditions, with exact receivers]
\label{thm:sharpcriterion}
The operations in Definition~\ref{def:sharp} make \(K^\sharp\) an object
of \(\SSS\) if and only if \(K\) is zerosumfree and has no nonzero
zero divisors.

For a standard semiring \(B\), a unital zero-preserving homomorphism
\(f:K\to B\), and \(b\in B\), consider the receiver equations
\[
 b+b=b,\quad b^2=b,\quad b+f(a)=b\ (a\in K),\quad
 bf(a)=b\ (a\in K\setminus\{0_K\}).
\tag{17}\label{eq:sharpreceivers}
\]
When the two stated conditions on \(K\) hold, restriction and evaluation
at \(w\) give a bijection between homomorphisms \(K^\sharp\to B\)
and pairs \((f,b)\) satisfying \eqref{eq:sharpreceivers}.
When either condition fails, every such unital receiver \(B\) is the
singleton, and the singleton is their universal receiver.
\end{theorem}
\begin{proof}
For necessity, if \(a,b\ne0_K\) and \(ab=0_K\), associativity in
the proposed extension would require
\[
 (wa)b=wb=w=w(ab)=w0_K=0_K,
\]
contradicting the disjointness of \(w\) and \(K\). If \(a+b=0_K\)
and at least one of \(a,b\) is nonzero, both must be nonzero: a zero
summand would leave the other unchanged. Distributivity would then require
\[
 0_K=w(a+b)=wa+wb=w+w=w,
\]
again a contradiction.

For sufficiency assume both conditions. Addition is commutative,
associative, and has identity \(0_K\): triples in \(K\) use its laws,
and any triple containing \(w\) has sum \(w\) under either
parenthesization. Multiplication is commutative by definition, and has
identity \(1_K\) because \(1_K\ne0_K\) and hence \(1_Kw=w\).
The element \(0_K\) absorbs every element, including \(w\).
For associativity, a triple with a zero factor gives zero under both
parenthesizations. A triple with no zero factor and no \(w\) uses
associativity in \(K\). A triple with no zero factor and an occurrence
of \(w\) gives \(w\) under both parenthesizations: a product of the
other nonzero elements of \(K\) cannot become zero, by hypothesis.

For distributivity \(a(b+c)=ab+ac\), first let \(a=0_K\); both sides
are \(0_K\). Next let \(a=w\). In the extended additive monoid,
\(b+c=0_K\) holds exactly when \(b=c=0_K\), by zerosumfreeness
on \(K\) and the rule for a sum with \(w\). If both summands are
zero, both sides of distributivity are zero. Otherwise the left side
is \(w\), while at least one of \(wb,wc\) is \(w\) and the other
is either zero or \(w\), so the right side is \(w\). Finally let
\(a\in K\setminus\{0_K\}\). If both \(b,c\) lie in \(K\), use
its original distributivity. If either is \(w\), their sum is \(w\)
and the left side is \(w\); the right side has a summand \(w\) and
therefore equals \(w\). These exhaust the cases and prove all axioms.

When the axioms hold, necessity of the receiver equations follows by
applying a homomorphism to the defining equations of the extension.
For sufficiency define the function to equal \(f\) on \(K\) and to
send \(w\) to \(b\). Pairs in \(K\) are handled by \(f\). An
additive pair with \(w\) is handled by the first or third equation
in \eqref{eq:sharpreceivers}. A multiplicative pair with \(w\) and
a nonzero element of \(K\), or with another \(w\), is handled by
the fourth or second equation. The remaining pair \((w,0_K)\) is
handled by the standard-semiring law \(b0_B=0_B=f(0_K)\).
The two constants are preserved by \(f\). This gives the unique extension.

Suppose now that \(a,b\in K\setminus\{0_K\}\) satisfy \(ab=0_K\),
and a receiver \((f,c)\) satisfies \eqref{eq:sharpreceivers}. Then
\[
 c=cf(b)=(cf(a))f(b)=cf(ab)=c0_B=0_B.
\]
If instead \(a+b=0_K\) with nonzero \(a,b\), then
\[
 0_B=cf(a+b)=cf(a)+cf(b)=c+c=c.
\]
In either case the equation \(c+f(1_K)=c\) gives
\(0_B+1_B=0_B\), so \(1_B=0_B\), and the standard-semiring laws
make \(B\) a singleton. The singleton has a unique map from \(K\)
and a unique choice of \(c\), satisfying every receiver equation.
Every receiver factors through it uniquely, which proves the universal
assertion even when the proposed distinct-element extension cannot exist.
\end{proof}

\begin{corollary}\label{cor:T}
For every commutative unital ring \(R\), the construction
\[
 T(R)=G(R)^\sharp=R\sqcup\{\tau,w\}
\tag{18}\label{eq:T}
\]
is a commutative unital standard semiring, with \(\tau\) as zero and
\(1_R\) as unit. Its inclusion \(G(R)\hookrightarrow T(R)\) is
an injective morphism in \(\SSS\). For \(r,s\in R\), its complete
stratum-level operations are
\[
\begin{array}{c|ccc}
 +&\tau&s&w\\\hline
 \tau&\tau&s&w\\
 r&r&r+s&w\\
 w&w&w&w
\end{array}
\qquad
\begin{array}{c|ccc}
 \cdot&\tau&s&w\\\hline
 \tau&\tau&\tau&\tau\\
 r&\tau&rs&w\\
 w&\tau&w&w
\end{array}.
\tag{19}\label{eq:Ttable}
\]
Here \(r,s\) range over \emph{all} elements of \(R\), including
\(e\), so in particular \(we=w\) but \(w\tau=\tau\).
These rules are the unique binary operations on this disjoint union
that preserve the operations on \(G(R)\), make \(w\) additively
absorbing, make \(\tau\) globally multiplicatively absorbing, and
make \(w\) multiply every element other than \(\tau\) to \(w\).
\end{corollary}
\begin{proof}
Proposition~\ref{prop:G} proves the exact two hypotheses of
Theorem~\ref{thm:sharpcriterion} relative to the zero \(\tau\).
Also \(1_R\ne\tau\) because the union is disjoint, even for the
zero ring. Thus the theorem applies with \(K=G(R)\). The operations
restricted to \(G(R)\), and both constants, are unchanged, proving
the inclusion assertion. Each table entry is either an original
\(G(R)\) operation, a zero-absorption rule, or one of the defining
\(w\) rules. Every pair in the disjoint union has exactly one such
specified value, proving the uniqueness assertion.
\end{proof}

\begin{proposition}[Exact difference between the two constructions]
\label{prop:difference}
The identification of underlying sets
\(F(R)\to T(R)\) that fixes \(G(R)\) and sends \(\Om\) to \(w\)
is not a multiplicative homomorphism, and neither is its inverse.
Every \(\HHZ\)-morphism \(F(R)\to T(R)\) is constant \(\tau\);
there is no \(\HH\)-morphism of that type.
If \(R\) is nonzero, there is no \(\HH\)-morphism \(T(R)\to F(R)\).

In \(T(R)\), the congruence generated by \(\tau=e\) is universal.
Every zero-preserving, possibly nonunital, map from \(T(R)\) to a
ring is constant zero. For nonzero \(R\), the projection
\(p_R:G(R)\to R\) therefore does not extend over \(T(R)\).
These conclusions concern \(T(R)\); for the literal object \(F(R)\),
the quotient by \(\tau=e\) is instead the exact object \(A(R)\)
in Theorem~\ref{thm:quotientA}.
\end{proposition}
\begin{proof}
The two multiplication tables disagree at \((\tau,\Om)\), respectively
\((\tau,w)\): the literal result is \(\Om\), while the alternate
result is \(\tau\). Since these are distinct, the indicated
identification and its inverse fail their multiplication equation.
The assertions about maps \(F(R)\to T(R)\) follow from
Theorem~\ref{thm:reflection}, because \(T(R)\) is a nontrivial
standard semiring: its distinct elements include \(\tau,e\).

Suppose \(R\) is nonzero and \(h:T(R)\to F(R)\) is an
\(\HH\)-morphism. Since \(w\tau=\tau\), its image satisfies
\(h(w)\tau=\tau\), so \(h(w)\ne\Om\) and \(h(w)\in G(R)\).
On the other hand \(w+1_R=w\) forces
\(h(w)+1_R=h(w)\). If \(h(w)=\tau\), this would give
\(1_R=\tau\), impossible. If \(h(w)=r\in R\), this is the
ring equation \(r+1_R=r\), implying \(1_R=e\), also impossible.

In a congruence of \(T(R)\) identifying \(\tau,e\), multiplication
by \(w\) identifies \(\tau w=\tau\) with \(ew=w\). Thus it
identifies \(\tau,w\). Adding any \(x\in T(R)\) identifies
\(\tau+x=x\) with \(w+x=w\), so all elements are equivalent.
For a binary-addition-preserving map to a ring, the equation
\(w+x=w\) and additive cancellation in the target force every image
to be zero, just as in the earlier ring-receiver proof. In particular
this holds for all of the stated homomorphisms. If \(R\) is nonzero,
the map \(p_R\) fixes \(1_R\ne e\), so cannot be the restriction of
such a constant-zero map.
\end{proof}

\begin{proposition}[Functorial scope of the alternate construction]
Let \(K,L\) be nontrivial commutative unital standard semirings meeting
the two conditions of Theorem~\ref{thm:sharpcriterion}. For a unital
zero-preserving map \(f:K\to L\), the set map
\[
 f^\sharp:K^\sharp\to L^\sharp,\qquad
 f^\sharp|_K=f,\quad f^\sharp(w_K)=w_L
\]
is a homomorphism if and only if \(f\) reflects zero:
\(f(a)=0_L\) implies \(a=0_K\).
Consequently every ring homomorphism \(h:R\to R'\) does extend to
a homomorphism \(T(R)\to T(R')\) fixing the two adjoined strata,
even if the ring map has a nonzero kernel.
\end{proposition}
\begin{proof}
If \(a\ne0_K\) but \(f(a)=0_L\), then
\(f^\sharp(w_Ka)=w_L\), while
\(f^\sharp(w_K)f^\sharp(a)=w_L0_L=0_L\), contradicting preservation
of multiplication. Conversely, if \(f\) reflects zero, pairs from
\(K\) obey the homomorphism equations for \(f\); additive pairs
with \(w_K\) map to the absorber equation; multiplicative pairs with
\(w_K\) and a nonzero entry map to \(w_L\) because that entry still
has nonzero image; the pair with zero maps to zero, and the pair of
two new elements maps to \(w_L^2=w_L\). Constants are preserved.
For a ring map \(h\), the map \(G(h):G(R)\to G(R')\) sends
\(\tau\) to \(\tau'\) and every \(r\in R\) into \(R'\), never
into the disjoint element \(\tau'\). It therefore reflects the
standard-semiring zero, regardless of whether \(h(r)=0_{R'}\)
for a nonzero ring element. Applying the proved criterion gives the
asserted map on \(T\).
\end{proof}

\end{document}
